\documentclass[11pt,a4paper]{article}

\usepackage{amsmath,amssymb}
\usepackage{physics}
\usepackage{booktabs}
\usepackage{geometry}
\usepackage{hyperref}
\usepackage{enumitem}

\geometry{margin=1in}
\hypersetup{colorlinks=true, linkcolor=blue, urlcolor=blue}

\title{Quantum Protein Folding Simulation Circuit \\[4pt]
       \large Explanation of \texttt{quantum\_folding\_circuit.ipynb}}
\author{}
\date{}

\begin{document}
\maketitle

\section{Overview}

This document explains the variational quantum circuit template implemented in
\texttt{quantum\_folding\_circuit.ipynb}. The notebook constructs a parameterised
ansatz for solving the protein folding problem on a quantum computer using the
Variational Quantum Eigensolver (VQE) framework, following the lattice model of
Robert et al.\ (2021).

%──────────────────────────────────────────────────────────────────────
\section{Tetrahedral Lattice Encoding}

Proteins are folded on a three-dimensional tetrahedral (diamond) lattice. Each
bond between consecutive amino acids is represented by one of four tetrahedral
bond vectors:

\begin{align}
  \vec{v}_0 &= \frac{1}{\sqrt{3}}(+1,\,+1,\,+1), &
  \vec{v}_1 &= \frac{1}{\sqrt{3}}(+1,\,-1,\,-1), \notag\\[4pt]
  \vec{v}_2 &= \frac{1}{\sqrt{3}}(-1,\,+1,\,-1), &
  \vec{v}_3 &= \frac{1}{\sqrt{3}}(-1,\,-1,\,+1).
\end{align}

The angle between any two distinct bond vectors is the tetrahedral angle:

\begin{equation}
  \theta_{\mathrm{tet}}
    = \arccos\!\left(-\frac{1}{3}\right)
    \approx 109.47^{\circ}.
\end{equation}

The lattice is bipartite (type-A / type-B sublattices), so consecutive residues
alternate between stepping by $+\vec{v}_i$ and $-\vec{v}_i$, preserving the
tetrahedral bond angle at every vertex.

%──────────────────────────────────────────────────────────────────────
\section{Dense Qubit Encoding}

Under the \emph{dense} encoding scheme, each turn direction is represented by
two qubits $(q_1, q_2)$. The chosen lattice axis is identified by four
\textbf{turn indicator functions}, exactly one of which evaluates to~1 for any
given qubit state:

\begin{alignat}{3}
  f_0 &= (1-q_1)(1-q_2) &\qquad& \ket{00} &\;\mapsto\;& \text{Axis } 0, \\
  f_1 &= q_2\,(q_2-q_1)  &\qquad& \ket{01} &\;\mapsto\;& \text{Axis } 1, \\
  f_2 &= q_1\,(q_1-q_2)  &\qquad& \ket{10} &\;\mapsto\;& \text{Axis } 2, \\
  f_3 &= q_1 \cdot q_2    &\qquad& \ket{11} &\;\mapsto\;& \text{Axis } 3.
\end{alignat}

An alternative \emph{sparse} encoding uses 4~qubits per turn in a one-hot
representation.

%──────────────────────────────────────────────────────────────────────
\section{Qubit Budget}

For a protein chain of $N$ amino acids there are $N-1$ bonds and therefore
$N-2$ turns. The first two turns are \textbf{pinned} to eliminate global
rotational and translational symmetry, leaving $N-3$ active turns.

\begin{equation}
  n = N_{\mathrm{cf}} + N_{\mathrm{in}},
\end{equation}

where the number of \textbf{configuration qubits} under dense encoding is

\begin{equation}
  N_{\mathrm{cf}} = 2(N-3),
\end{equation}

and $N_{\mathrm{in}}$ is the number of auxiliary \textbf{interaction qubits}
used to track physical contacts between non-adjacent residues.

\paragraph{Example (7-mer peptide).}
\begin{equation}
  N_{\mathrm{cf}} = 2(7-3) = 8, \qquad
  N_{\mathrm{in}} = 1, \qquad
  n = 9.
\end{equation}

%──────────────────────────────────────────────────────────────────────
\section{Variational Ansatz Circuit \texorpdfstring{$U(\vec\theta)$}{U(theta)}}

The parameterised circuit is composed of three stages.

\subsection{Step A --- Initialisation Layer}

A Hadamard gate creates an equal superposition on every qubit, followed by a
layer of parameterised $R_Y$ rotations:

\begin{equation}
  U_{\mathrm{init}}
    = \prod_{i=0}^{n-1} R_Y(\theta_i)\; H_i\,.
\end{equation}

\subsection{Step B --- Entangling and Rotation Blocks}

The following block is repeated $m$ times (where $m = \texttt{depth}$). For
each layer $d = 1, \dots, m$:

\begin{equation}
  U_d
    = \left[\,\prod_{i=0}^{n-1} R_Y(\theta_{d,i})\right]
      \cdot U_{\mathrm{ent}}\,.
\end{equation}

\paragraph{Closed-loop entangling layer.}
In the default \texttt{closed\_loop} mode the entangling unitary is a linear
CNOT chain with a loopback from the last qubit to the first:

\begin{equation}
  U_{\mathrm{ent}}
    = \mathrm{CNOT}(n{-}1,\,0)
      \;\cdot\; \prod_{i=0}^{n-2} \mathrm{CNOT}(i,\,i{+}1)\,.
\end{equation}

This topology is well-suited to NISQ hardware with nearest-neighbour
connectivity.

\paragraph{All-to-all entangling layer.}
The alternative \texttt{all\_to\_all} mode applies a CNOT between every pair of
qubits:

\begin{equation}
  U_{\mathrm{ent}}
    = \prod_{i=0}^{n-1}\;\prod_{j=i+1}^{n-1}
      \mathrm{CNOT}(i,\,j)\,.
\end{equation}

\subsection{Step C --- Measurement}

All $n$ qubits are measured in the computational ($Z$) basis.

%──────────────────────────────────────────────────────────────────────
\section{Parameter Count}

The total number of variational parameters is

\begin{equation}
  \lvert\vec\theta\rvert = n\,(m+1)\,.
\end{equation}

For the 7-mer example ($n=9$, $m=2$):

\begin{equation}
  \lvert\vec\theta\rvert = 9 \times 3 = 27.
\end{equation}

%──────────────────────────────────────────────────────────────────────
\section{Gate Count Summary (7-mer Example)}

\begin{table}[h]
\centering
\begin{tabular}{@{}lc@{}}
  \toprule
  \textbf{Gate}        & \textbf{Count} \\
  \midrule
  $H$                  & 9  \\
  $R_Y(\theta)$        & 27 \\
  CNOT                 & 18 \\
  Measure              & 9  \\
  \midrule
  Circuit depth        & 23 \\
  \bottomrule
\end{tabular}
\caption{Gate counts for the 7-amino-acid ansatz with depth $m=2$
         and closed-loop entanglement.}
\end{table}

%──────────────────────────────────────────────────────────────────────
\section{Classical Post-Processing}

Given a measured bitstring $\mathbf{b} \in \{0,1\}^n$, the turn direction for
turn~$t$ is decoded by extracting $(q_1, q_2)$ from the configuration register
and evaluating the indicator functions $f_0, \dots, f_3$. Exactly one indicator
equals~1, identifying the chosen lattice axis $\alpha(t) \in \{0,1,2,3\}$.

The three-dimensional position of residue $k$ is then reconstructed as

\begin{equation}
  \vec{r}_k
    = \vec{r}_0
      + \sum_{t=1}^{k-1} (-1)^{t+1}\, \vec{v}_{\,\alpha(t)}\,,
\end{equation}

where the alternating sign $(-1)^{t+1}$ enforces the bipartite sublattice
structure of the tetrahedral lattice.

%──────────────────────────────────────────────────────────────────────
\section{VQE Context}

The ansatz $U(\vec\theta)$ produces candidate folding conformations as quantum
states. A classical optimiser tunes $\vec\theta$ to minimise the expectation
value of the protein's energy Hamiltonian:

\begin{equation}
  \vec\theta^{*}
    = \arg\min_{\vec\theta}\;
      \bra{0} U^{\dagger}(\vec\theta)\, H_{\mathrm{fold}}\, U(\vec\theta) \ket{0}\,.
\end{equation}

The Hamiltonian $H_{\mathrm{fold}}$ (not constructed in the notebook) typically
includes nearest-neighbour interaction energies and penalty terms that enforce
the self-avoidance constraint on the lattice.

\end{document}
